\begin{proposition}[The complete graded anticanonical ring]\hypertarget{tag-01S6}{}\label{inv:B2.2}
\StatementRecord{inv:B2.2}{proved}{01S6}{\SeriesRef{inv:B2.1}}

Put $m_1=3k$, $m_2=4l$, and $d=4l+3k-kl$.
Choose integers $r,s$ with $m_2r+m_1s=1$ and set
$\mathcal L=\OO_X(rS_1+sS_2)$.
There is a graded algebra isomorphism
\[
\bigoplus_{n\ge0}H^0(X,-nK_X)
\simeq\bigoplus_{n\ge0}\C[x,y]_{dn},
\qquad\deg x=m_2,\quad\deg y=m_1,
\]
where the right-hand summand is assigned degree $n$.
For negative weights the polynomial space is zero. Explicitly
\[
h^0(X,-nK_X)=
\max\left\{1+n-\left\lceil\frac{n(k-1)}{3k}\right\rceil
-\left\lceil\frac{n(3l-1)}{4l}\right\rceil,\,0\right\}.
\]
The integer $d$ is nonzero for admissible pairs. If $d<0$ the
anticanonical ring is $\C$; if $d>0$ it is the indicated weighted
Veronese ring. For $(k,l)=(1,1)$ it is
$\C[U,V]$, $\deg U=1$, $\deg V=2$.
\end{proposition}


\begin{proof}
\begin{enumerate}[label=\arabic*.,leftmargin=*]
\item The divisor of $f^*(t/(t-1))$ is $m_1S_1-m_2S_2$.
It follows, retaining this rational function as the comparison, that
$\mathcal L^{m_2}\simeq\OO(S_1)$,
$\mathcal L^{m_1}\simeq\OO(S_2)$ and
$\mathcal L^{m_1m_2}\simeq f^*\OO(1)$.
\ref{inv:B2.1} then gives $-K_X=\mathcal L^d$.

\item For every integer
power of $\mathcal L$, a meromorphic function representing a section
of $\OO(aS_1+bS_2)$ is holomorphic on every regular compact fibre
and is constant there. It descends to the base on the smooth locus.
On a cyclic cover $q=s^m$, the finite pole bound gives a Laurent
expansion bounded below, and invariance leaves only powers divisible
by $m$. Thus descent extends meromorphically across the finite point,
with the exact valuation inequality $m\operatorname{ord}_q(h)+a\ge0$.
The identical argument at the reduced cusp has $m=1$.
Conversely each base function satisfying these bounds pulls back to
an allowed section. Consequently, as subsheaves of meromorphic functions,
\[
f_*\OO(aS_1+bS_2)=
\OO_B\left(\left\lfloor a/m_1\right\rfloor p_1+
\left\lfloor b/m_2\right\rfloor p_2\right).
\]
This proves the displayed dimension formula using \ref{inv:B2.1} and
$h^0(\Pone,\OO(e))=\max\{e+1,0\}$.

\item For the algebra, let $x,y$ be the canonical sections of
$\OO(S_1),\OO(S_2)$ under the chosen identifications with powers of
$\mathcal L$. Their powers $x^{m_1},y^{m_2}$ pull back two independent
sections of $\OO_B(1)$, with zeros at the two distinct points.
For any integer $j$, choose the unique residues
$0\le a_0<m_1$, $0\le b_0<m_2$ for which
$j-m_2a_0-m_1b_0=e m_1m_2$ is integral.
The valuation formula identifies $H^0(\mathcal L^j)$ with
$x^{a_0}y^{b_0}H^0(\OO_B(e))$. When $e\ge0$ a basis is
\[
x^{a_0+m_1i}y^{b_0+m_2(e-i)},\qquad0\le i\le e;
\]
when $e<0$ there are no sections. These are exactly the monomials of
weighted degree $j$. Their multiplication is ordinary polynomial
multiplication, proving the algebra assertion in every degree.
When $k=l=1$, the weight-$6n$ monomials are generated by $y^2$ and
$x^3$, of regraded degrees $1,2$.
Finally, $d=0$ would imply $k\mid4l$ and $l\mid3k$; coprimality
forces $k=l=1$, which instead gives $d=6$.
\end{enumerate}
\end{proof}